\documentclass{article}

\newif\ifanonsubmission
\anonsubmissiontrue
\newif\ifneuripsstyle
\newif\ifincludeappendix
\newif\ifincludechecklist
\includeappendixfalse
\includechecklisttrue

\IfFileExists{neurips_2026.sty}{
  \neuripsstyletrue
  \ifanonsubmission
    \usepackage[eandd]{neurips_2026}
  \else
    \usepackage[eandd,preprint]{neurips_2026}
  \fi
}{
  \IfFileExists{neurips_2025.sty}{
    \neuripsstyletrue
    \ifanonsubmission
      \usepackage{neurips_2025}
    \else
      \usepackage[preprint]{neurips_2025}
    \fi
  }{
    \neuripsstylefalse
    \usepackage[margin=1in]{geometry}
  }
}

\usepackage[T1]{fontenc}
\usepackage[utf8]{inputenc}
\usepackage{microtype}
\usepackage{amsmath,amssymb}
\usepackage{booktabs}
\usepackage{graphicx}
\usepackage{xcolor}
\usepackage{enumitem}
\usepackage{multirow}
\usepackage{float}
\usepackage[section]{placeins}
\usepackage{url}
\usepackage{hyperref}
\makeatletter
\setlength{\@fptop}{0pt}
\setlength{\@fpbot}{0pt plus 1fil}
\setlength{\@dblfptop}{0pt}
\setlength{\@dblfpbot}{0pt plus 1fil}
\makeatother
\ifneuripsstyle\else
  \usepackage[numbers]{natbib}
\fi

\providecommand{\answerYes}[1][]{\textbf{[Yes]}#1}
\providecommand{\answerNo}[1][]{\textbf{[No]}#1}
\providecommand{\answerNA}[1][]{\textbf{[N/A]}#1}

\newcommand{\figdir}{../../outputs/paper-freeze-2026-04-18/figures}

\title{LatentGoalOps: A Benchmark for Latent Objective Adaptation Under Silent Shifts}
\ifanonsubmission
  \author{Anonymous Author(s)}
\else
  \author{Anonymous Author(s)}
\fi
\date{}

\begin{document}
\maketitle

\begin{abstract}
Latent objectives are common in real operational settings: stakeholders reveal
constraints, preferences, and feedback over time, while the true utility
function governing good decisions may remain implicit or change without
announcement. We introduce \textbf{LatentGoalOps}, a benchmark for evaluating
latent objective inference and adaptation under silent objective shifts. Agents
observe a fixed public operating schema consisting of business state,
stakeholder messages, decision ledgers, delayed effects, and governance
constraints, but never observe the hidden objective used by the deterministic
grader. The benchmark contains three sequential business-operations
environments; our main analysis focuses on two week-long anchor tasks, with a
shorter planning task used as breadth evidence. Across the anchor tasks,
restricted-view leakage audits show no easy shortcut, models incur a persistent
penalty on silently shifted episodes, and the best current local model remains
roughly 0.20 below a myopic oracle that observes the active hidden objective
but does not solve the full future trajectory. We further find that performance
is mechanism-sensitive: several sharp heldout drops are caused by
structured-action failures rather than latent-policy quality alone. For this
reason, LatentGoalOps reports strict score, native-only score, and interface
reliability jointly. These results position LatentGoalOps as a benchmark for
studying hidden-objective adaptation under a fixed public schema and
deterministic grader, rather than as a measure of general business-agent
readiness. Small hosted-model probes show heterogeneous interface reliability,
but none closes the oracle gap.
\end{abstract}

\section{Introduction}

Many agent benchmarks make the objective explicit: the instruction is written
down, the tool contract is visible, and success is defined against an external
goal. Recent suites such as \textsc{$\tau$}-bench and BrowserGym have pushed
that setting toward longer-horizon, tool-using evaluation in realistic
interfaces~\citep{taubench,browsergym}. Real operating environments, however,
often require a different capability. Leadership priorities are not always
fully stated, stakeholder messages are noisy, visible KPI movement is
ambiguous, and the underlying objective can change without announcement.

This paper targets that gap. We study \emph{latent objective inference and
adaptation}: inferring an unstated reward structure from public evidence,
acting coherently through a strict interface, and adapting when the hidden
objective changes mid-episode. We instantiate this construct in
\textbf{LatentGoalOps}, a synthetic but structured operating-world benchmark
with persistent entities, delayed effects, deterministic grading, and silent
objective shifts. The synthetic construction is intentional: it gives the
benchmark hidden-objective ground truth while preserving a realistic operating
surface for the agent.

The release contains three paper-facing sequential environments:
\emph{P1 Startup Week}, \emph{P2 Incident Response Week}, and
\emph{P3 Quarterly Headcount Plan}. The paper centers P1 and P2 as anchor
tasks; P3 serves only as shorter breadth evidence. Within that scope, the
frozen evidence supports three claims: current local open-weight
configurations retain a clear myopic-oracle gap, within-family behavior is
heterogeneous rather than uniformly monotone, and some large heldout drops are
interface-mediated rather than pure policy regressions.

\paragraph{Related work and positioning.}
Existing interactive benchmarks such as \textsc{$\tau$}-bench and BrowserGym
primarily evaluate explicit user goals in realistic interfaces, while
structured-output and function-calling evaluations such as
JSONSchemaBench and BFCL isolate interface control itself
~\citep{taubench,browsergym,jsonschemabench,bfcl}. LatentGoalOps instead hides
the objective, introduces delayed effects and silent shifts, and requires
strict structured actions under non-stationarity. Following recent
benchmark-methodology work, we therefore report strict score, native-only
score, and interface reliability together rather than a single leaderboard
scalar~\citep{agentrewardbench,betterbench,dataperf,benchmarkrepos}.

\begin{table}[t]
  \centering
  \scriptsize
  \setlength{\tabcolsep}{4pt}
  \begin{tabular}{lccccc}
    \toprule
    Benchmark & Business ops & Hidden objective & Silent shift & Delayed effects & Strict schema \\
    \midrule
    WebArena & No & No & No & No & No \\
    WorkArena & Yes & No & No & No & No \\
    BrowserGym & No & No & No & No & No \\
    OSWorld & No & No & No & No & No \\
    \textsc{$\tau$}-bench & Yes & No & No & Limited & No \\
    LatentGoalOps & Yes & Yes & Yes & Yes & Yes \\
    \bottomrule
  \end{tabular}
  \caption{Qualitative construct comparison against nearby interactive-agent
  benchmarks~\citep{webarena,workarena,browsergym,osworld,taubench}. LatentGoalOps
  differs by combining hidden objective, silent shift, delayed effects, and a
  strict executable action shell.}
  \label{tab:comparison}
\end{table}

\paragraph{Contributions.}
\begin{itemize}[leftmargin=1.5em]
  \item We introduce LatentGoalOps, a benchmark suite for latent objective
  inference and adaptation under silent non-stationarity with delayed effects,
  hidden strategic objectives, deterministic grading, two week-long anchor
  environments, and one shorter breadth task.
  \item We release a runnable benchmark artifact with a frozen protocol,
  core and heldout splits, deterministic baselines, restricted-view leakage
  probes, generated paper assets, and reviewer-facing reproducibility
  documentation.
  \item We report a strengthened local benchmark matrix with 630 primary model
  episodes and 360 baseline episodes, plus a P1/P2 anchor-task robustness
  slice that isolates the benchmark's strongest latent-adaptation settings.
  \item We add exploratory hosted-model probes for Gemini 3 Flash preview,
  GPT-5.4 mini, and Claude Sonnet 4.6 to test whether accessible API models
  immediately close the myopic-oracle gap under the same action shell.
  \item We show that the benchmark's empirical story is mechanism-sensitive:
  some larger-model changes are native-policy movements, while some of the
  sharpest heldout regressions are interface-mediated drops.
\end{itemize}

\section{Benchmark}

\subsection{Environment Loop}

Each episode is a finite partially observed decision process. At step $t$, the
agent receives a public observation $o_t$, emits a structured action $a_t$, and
the environment advances to the next state. Public observations include inbox
messages, KPI dashboards, dated decision ledgers, pending effects, realized
effects, customer accounts, stakeholder personas, team state, market context,
and visible governance constraints. The hidden objective $g_t$ is never shown
directly.

Formally, an episode is a trajectory
\[
\tau = (o_0, a_0, o_1, a_1, \ldots, o_T),
\]
with hidden state $x_t$, public observation $o_t = \psi(x_t)$, and active
hidden objective state $g_t$. The environment maintains a normalized
four-channel metric vector $m_t \in [0,1]^4$ over growth, retention, revenue,
and efficiency, and an active latent weight vector $w_t \in \Delta^4$. The
latent utility used by the grader is
\[
U_t = \mathrm{clip}_{[0,1]} \left( \sum_{c \in \mathcal{C}} m_{t,c} \, w_{t,c} \right),
\qquad
\mathcal{C} = \{\text{growth}, \text{retention}, \text{revenue}, \text{efficiency}\}.
\]

The action interface is intentionally strict. On the sequential tasks, a valid
action is a JSON-serializable payload over task-specific fields such as
\texttt{chosen\_initiatives}, \texttt{messaging\_action},
\texttt{pricing\_change\_pct}, \texttt{support\_policy}, or
\texttt{budget\_allocations}. If the model fails to emit a valid structured
action, the benchmark executes a fallback action and the trajectory still
advances. In other words, interface failures are part of the task, not merely
parser trivia. This follows the broader lesson from structured-generation and
tool-use evaluation that output control is part of agent capability rather than
an implementation footnote~\citep{jsonschemabench,bfcl}.

\begin{figure}[t]
  \centering
  \includegraphics[width=\linewidth]{\figdir/benchmark_overview.pdf}
  \caption{LatentGoalOps exposes public operating state and delayed
  consequences while hiding the objective that determines utility. The agent
  must act through a strict structured-action interface; failures at that
  interface alter the trajectory itself.}
  \label{fig:overview}
\end{figure}

\subsection{Three Sequential Environments and Silent Shifts}

LatentGoalOps evaluates three sequential operating environments. The
implementation preserves internal task ids for reproducibility, but the main
paper uses the cleaner P1--P3 names throughout. Table~\ref{tab:tasksuite}
summarizes the benchmark environments and their roles in the manuscript.

\begin{table}[t]
  \centering
  \scriptsize
  \setlength{\tabcolsep}{4pt}
  \begin{tabular}{l l c c c}
    \toprule
    Paper ID & Environment & Horizon & Paper role & Hidden shift \\
    \midrule
    P1 & Startup Week & 10 & anchor & Yes \\
    P2 & Incident Response Week & 7 & anchor & Yes \\
    P3 & Quarterly Headcount Plan & 4 & breadth & Yes \\
    \bottomrule
  \end{tabular}
  \caption{LatentGoalOps benchmark environments. The main paper's
  result-bearing analyses center the two week-long anchor tasks P1 and P2.}
  \label{tab:tasksuite}
\end{table}

P1 Startup Week is a ten-step operating week with initiatives, renewals, and
delayed business consequences. P2 Incident Response Week is a seven-step
incident setting where trust repair, governance, recovery speed, and cost can
conflict. P3 Quarterly Headcount Plan is a shorter staffing task used only for
breadth. Heldout evaluation changes scenario families and delayed-effect
realizations while preserving the same public schema and grading semantics.
Silent shifts occur with probability $0.6$, usually at steps $\{3,4,5\}$, and
can alter latent weights as well as risk posture, planning horizon, segment
focus, and governance strictness. The model is never told that the shift
occurred.

\subsection{Scoring and Reporting}

For the primary sequential tasks, the episode score is an explicit weighted
combination of final latent utility $U$, adaptation $A$, coherence $C$, and
constraint satisfaction $K$:
\[
S_{\mathrm{task}} = \alpha_U U + \alpha_A A + \alpha_C C + \alpha_K K.
\]
Task-specific coefficients and no-shift renormalization rules are fixed in the
release protocol. This keeps the main paper focused on
the measurement construct rather than coefficient bookkeeping.

For \emph{Startup Week} and \emph{Incident Response Week}, adaptation measures
how quickly post-shift decisions improve under the active hidden objective and
whether that recovery is sustained; P3 uses the same principle with a shorter
horizon. We report overall strict score, native-only score, shifted versus
unshifted score, and interface reliability. Release aggregates pool 10 seeds
per task across P1--P3, while P1/P2 anchor slices isolate the strongest
construct signal; aggregate views use 95\% bootstrap confidence intervals.

\section{Evaluation and Benchmark Validity}

\subsection{Baseline Separation}

LatentGoalOps has meaningful score dynamic range on the frozen release
aggregate over P1--P3. Random, visible-heuristic, and myopic-oracle baselines
preserve the expected ordering on both splits, and the myopic-oracle gap
remains substantial on both splits. The myopic oracle
observes the active hidden objective and chooses the best immediate action, but
does not solve the full future trajectory; we use it as a construct headroom
reference rather than as a fair agent baseline or human proxy.

\begin{table}[t]
  \centering
  \small
  \begin{tabular}{llrrr}
    \toprule
    Baseline & Split & Episodes & Mean score & 95\% CI \\
    \midrule
    random & core & 60 & 0.303 & [0.286, 0.320] \\
    visible heuristic & core & 60 & 0.346 & [0.329, 0.362] \\
    myopic oracle & core & 60 & 0.599 & [0.575, 0.623] \\
    random & heldout & 60 & 0.310 & [0.290, 0.330] \\
    visible heuristic & heldout & 60 & 0.371 & [0.356, 0.388] \\
    myopic oracle & heldout & 60 & 0.598 & [0.575, 0.623] \\
    \bottomrule
  \end{tabular}
  \caption{Baseline separation on the frozen primary-task aggregate.}
  \label{tab:baselines}
\end{table}


\subsection{Leakage Audit}

Restricted-view centroid probes test for shortcut access to the hidden
objective. Table~\ref{tab:leakagecompact} shows that P1 and P2 stay near the
four-way chance regime under both full and restricted views, with
\texttt{explicit\_goal\_token\_rate} equal to 0.0 throughout. P3 is less
clean: its core restricted-view maximum reaches 0.35, which is why we treat it
as breadth evidence rather than as a validity anchor. Each probe slice uses 20
train and 20 test examples, so this audit is a shortcut screen rather than a
proof of impossibility.

\begin{table}[t]
  \centering
  \scriptsize
  \setlength{\tabcolsep}{4pt}
  \begin{tabular}{llrrrrrr}
    \toprule
    Task & Split & Full & Numeric & Sender & Order & Max restricted & Goal tokens \\
    \midrule
    Startup Week & core & 0.20 & 0.05 & 0.25 & 0.25 & 0.30 & 0.00 \\
    Startup Week & heldout & 0.20 & 0.15 & 0.20 & 0.25 & 0.25 & 0.00 \\
    Incident Response Week & core & 0.20 & 0.10 & 0.30 & 0.25 & 0.30 & 0.00 \\
    Incident Response Week & heldout & 0.15 & 0.20 & 0.30 & 0.25 & 0.30 & 0.00 \\
    Quarterly Headcount Plan & core & 0.25 & 0.25 & 0.25 & 0.25 & 0.35 & 0.00 \\
    Quarterly Headcount Plan & heldout & 0.20 & 0.15 & 0.25 & 0.25 & 0.30 & 0.00 \\
    \bottomrule
  \end{tabular}
  \caption{Compact restricted-view leakage summary for paper-facing tasks.
  Chance is 0.25 for the four-way archetype probe. ``Max restricted'' is the
  maximum over non-full views in the frozen audit and is interpreted as a
  screening diagnostic rather than a proof of no shortcut.}
  \label{tab:leakagecompact}
\end{table}


Utility-only rescoring sharply compresses the myopic-oracle gap, which is the
expected pattern if the benchmark is measuring latent-objective adaptation
rather than only final-state utility.

\section{Main Results}

\subsection{The Two Week-Long Anchor Tasks Carry the Main Result}

Because P3 is shorter-horizon and less discriminative than the two week-long
environments, Table~\ref{tab:anchorslice} reports a P1/P2-only robustness
slice computed directly from the frozen release as the paper's main result
view. On the anchor slice, the strongest local core mean is Qwen
27B at 0.425 and the strongest heldout mean is DeepSeek 8B at 0.431, while the
myopic-oracle anchor means are 0.626 on core and 0.630 on heldout. The best
local anchor results therefore still trail the oracle by about 0.20.

\begin{table}[t]
  \centering
  \scriptsize
  \setlength{\tabcolsep}{4pt}
  \begin{tabular}{lrrrr}
    \toprule
    Model & Core P1/P2 & Heldout P1/P2 & Core native rate & Heldout native rate \\
    \midrule
    Myopic oracle & 0.626 & 0.630 & 1.00 & 1.00 \\
    \midrule
    DeepSeek 1.5B & 0.407 & 0.407 & 1.00 & 1.00 \\
    DeepSeek 8B & 0.408 & 0.431 & 1.00 & 1.00 \\
    DeepSeek 14B & 0.420 & 0.405 & 1.00 & 1.00 \\
    Gemma 2B & 0.401 & 0.412 & 1.00 & 1.00 \\
    Gemma 4B & 0.408 & 0.416 & 1.00 & 1.00 \\
    Gemma 26B & 0.392 & 0.341 & 0.65 & 0.70 \\
    GPT-OSS 20B & 0.407 & 0.358 & 1.00 & 1.00 \\
    Qwen 2B & 0.406 & 0.412 & 1.00 & 1.00 \\
    Qwen 9B & 0.394 & 0.393 & 1.00 & 1.00 \\
    Qwen 27B & 0.425 & 0.340 & 0.95 & 0.25 \\
    \bottomrule
  \end{tabular}
  \caption{P1/P2 robustness slice from the frozen release. The myopic-oracle
  row makes the remaining adaptation headroom explicit, while the weakest
  heldout cells coincide with shell-fragile models.}
  \label{tab:anchorslice}
\end{table}


Taken together with the leakage audit and the persistent shifted-versus-
unshifted gap, the P1/P2 anchor tasks make the intended construct visible:
public trajectory evidence is informative but incomplete, so successful
behavior requires latent-goal tracking rather than shortcut extraction.

\subsection{The Full Release Aggregate Preserves the Same Story}

The full P1--P3 aggregate preserves the same story: the strongest frozen local
aggregates are 0.443 on core and 0.446 on heldout, still about 0.15 below the
myopic-oracle baselines of 0.599 and 0.598.

\subsection{Hosted API Probes Improve Reliability but Do Not Close the Gap}

To test whether readily accessible hosted models close the benchmark's oracle
gap, we ran small OpenRouter probes on the P1--P3 aggregate after freezing the primary
local-model matrix. These runs are exploratory and not used as load-bearing
within-family scaling evidence, but they answer a practical reviewer question:
whether modern API models immediately erase the oracle gap. Because P1 and P2
are the load-bearing anchor tasks, Table~\ref{tab:apiprobes} reports both
P1/P2 means and the broader P1--P3 aggregate.

\begin{table}[t]
  \centering
  \scriptsize
  \setlength{\tabcolsep}{3pt}
  \resizebox{\linewidth}{!}{%
  \begin{tabular}{llrrrrrrr}
    \toprule
    Model & Split & Episodes & P1/P2 & P1--P3 & Native-only & Native n & Native rate & Empty fallback \\
    \midrule
    Gemini 3 Flash preview & core & 30 & 0.346 & 0.328 & 0.385 & 8 & 0.27 & 0.73 \\
    Gemini 3 Flash preview & heldout & 30 & 0.319 & 0.313 & 0.338 & 6 & 0.20 & 0.80 \\
    GPT-5.4 mini & core & 30 & 0.388 & 0.418 & 0.414 & 27 & 0.90 & 0.10 \\
    GPT-5.4 mini & heldout & 30 & 0.352 & 0.381 & 0.402 & 21 & 0.70 & 0.30 \\
    Claude Sonnet 4.6 & core & 30 & 0.423 & 0.429 & 0.429 & 30 & 1.00 & 0.00 \\
    Claude Sonnet 4.6 & heldout & 30 & 0.385 & 0.393 & 0.393 & 30 & 1.00 & 0.00 \\
    \bottomrule
  \end{tabular}
  }
  \caption{Exploratory hosted-model probes run through OpenRouter on the P1--P3
  aggregate. These API probes are excluded from the primary local-model matrix
  and are used only to test whether readily accessible hosted models close the
  gap to the myopic oracle. Entries are descriptive point estimates.}
  \label{tab:apiprobes}
\end{table}


The probes separate interface reliability from latent-objective quality.
Gemini 3 Flash preview often produced plausible intent with non-native field
names, yielding low native rates and many empty fallbacks. GPT-5.4 mini was
mostly native and reached 0.388/0.352 on the P1/P2 core/heldout anchor slice
(0.418/0.381 on P1--P3). Claude Sonnet 4.6 was the cleanest probe, with native
rate 1.00 on both splits and P1/P2 scores of 0.423/0.385. None of these
descriptive point estimates closes the myopic-oracle gap. Under these small
probes, better shell reliability does not by itself close the gap to a
privileged policy that sees the active hidden objective.

\subsection{Some of the Sharpest Heldout Drops Are Interface-Mediated}

Across the frozen release, scaling is family-specific rather than monotone, but
the most important mechanism is interface confounding rather than raw ranking
alone. DeepSeek remains the cleanest native-policy signal: its three reported
sizes stay at native rate 1.0 on both splits, so its gains and losses are best
read as policy movement. Qwen 2B $\rightarrow$ 9B also changes at native rate
1.0, giving a clean example of non-monotone native-policy behavior. The sharp
heldout failures, by contrast, are dominated by shell fragility.

\begin{figure}[t]
  \centering
  \includegraphics[width=0.9\linewidth]{\figdir/overall_vs_native.pdf}
  \caption{Observed score versus native-only score on the frozen release
  aggregate. Small-native-count estimates are diagnostic only, but the figure
  makes shell confounding visible instead of hiding it inside a single scalar
  mean.}
  \label{fig:overall-native}
\end{figure}

The clearest examples are Gemma 26B and Qwen 27B. On the P1/P2 anchor slice,
Gemma 26B heldout falls to 0.341 overall at native rate 0.70, while Qwen 27B
heldout falls to 0.340 overall at native rate 0.25. On the full release
aggregate, Qwen 27B heldout rises from 0.385 overall to 0.457 on the native
subset, but that subset contains only 14 episodes and must be read
diagnostically rather than as a leaderboard claim. Table~\ref{tab:qualtraj}
grounds that mechanism in exact stored outputs: DeepSeek stays incident-
focused across the shift, whereas Qwen 27B hits an empty fallback and later
returns stale growth/revenue moves.

\begin{table}[t]
  \centering
  \scriptsize
  \setlength{\tabcolsep}{3pt}
  \begin{tabular}{p{0.14\linewidth} p{0.18\linewidth} p{0.43\linewidth} p{0.15\linewidth}}
    \toprule
    Run / step & Public clue & Selected output fields & Strict / reward \\
    \midrule
    DeepSeek 14B, P2 day 2 & Day 2; 3 pending effects & {\ttfamily initiative=ship\_usage\_pricing; message=retention\_campaign; pricing=-0.1; support=incident\_swarm} & strict; $-0.0302$ \\
    DeepSeek 14B, P2 day 8 & Day 8; 0 pending effects & {\ttfamily initiative=refactor\_incident\_tooling; message=retention\_campaign; pricing=0.0; support=incident\_swarm} & strict; $0.1427$ \\
    Qwen 27B, P1 day 5 & Day 5; 0 pending effects & {\ttfamily \{\}} & non-strict; $-0.0216$ \\
    Qwen 27B, P1 day 9 & Day 9; 3 pending effects & {\ttfamily initiative=launch\_referral\_loop; message=growth\_push; pricing=0.0; support=balanced\_triage} & strict; $-0.2013$ \\
    \bottomrule
  \end{tabular}
  \caption{Exact selected fields copied from stored action payloads and episode
  logs in the frozen case studies. These are action outputs and scalar episode
  diagnostics, not chain-of-thought traces.}
  \label{tab:qualtraj}
\end{table}

\subsection{Silent Shifts Remain the Dominant Stressor}

Silent objective shifts remain the primary benchmark stress test. Across the
release matrix, shifted episodes are materially harder than unshifted ones on
both core and heldout, with a median unshifted-minus-shifted gap of 0.172 in
the release-wide diagnostics. This gap is not explained away by shell failures:
models with perfect native-action reliability, including all DeepSeek sizes,
still pay a substantial shift penalty. That is the benchmark construct we most
want to isolate. Figure~\ref{fig:shifted} makes this visible at the aggregate
level.

\begin{figure}[t]
  \centering
  \includegraphics[width=0.82\linewidth]{\figdir/shifted_vs_unshifted.pdf}
  \caption{Shifted episodes remain materially harder than no-shift episodes on
  the frozen primary-task aggregate. Error bars denote 95\% bootstrap
  confidence intervals for the score aggregates.}
  \label{fig:shifted}
\end{figure}

\section{Discussion}

LatentGoalOps clears the benchmark-science checks we care about most: baseline
ordering, a persistent oracle gap, a consistent shift penalty, and no easy
shortcut on the two anchor tasks. Because malformed actions change the
trajectory itself, strict score, native-only score, and native episode rate
need to be reported together.

\section{Limitations and Intended Use}

\begin{itemize}[leftmargin=1.5em]
  \item The benchmark is synthetic and intended for comparative evaluation, not
  operational advice or direct real-world deployment claims.
  \item The heldout split preserves schema and grader semantics, and P3 is
  treated as shorter planning breadth rather than as a load-bearing anchor.
  \item GPT-OSS 120B is API-served, and this version does not include human
  baselines or a full human-audited grader validation study.
  \item The hosted-model probes are intentionally small because of API access,
  budget, and deadline constraints. They should be read as exploratory
  external probes, not exhaustive frontier-model coverage.
  \item The adaptation metric uses thresholded detection rules; we expose and
  audit those choices in the release protocol, but smoother alternatives remain
  possible.
\end{itemize}

\section{Conclusion}

LatentGoalOps isolates hidden-objective adaptation under silent
non-stationarity. On the tighter P1/P2 anchor slice and the broader release
aggregate, the benchmark shows no easy shortcut on the anchor tasks, a
persistent shift penalty, and a clear gap between current local models and the
myopic-oracle reference, while also revealing that some sharp heldout failures
are interface-mediated. That combination argues for evaluating latent policy
quality and action-shell reliability separately.

\bibliographystyle{plainnat}
\bibliography{references}

\ifincludeappendix
  \clearpage
  \appendix
  \section{Protocol Snapshot}

This supplement is intentionally narrower than the main paper. It records the
frozen protocol, serving path, supplementary diagnostics, and provenance for
the strengthened bundle at \texttt{outputs/paper-freeze-2026-04-18}. It does
not repeat the main model table or baseline table, which are already reported
in the main paper. The paper-facing benchmark has three sequential
environments; the main manuscript centers P1 and P2 as the result-bearing
anchor tasks and treats P3 as shorter breadth evidence. Auxiliary one-shot
tasks in the code release are for smoke tests and extensibility, not headline
claims.

\begin{table}[!h]
  \centering
  \scriptsize
  \setlength{\tabcolsep}{4pt}
  \begin{tabular}{lllp{0.12\linewidth}p{0.38\linewidth}}
    \toprule
    Paper ID & Environment & Internal id & Horizon & Purpose \\
    \midrule
    P1 & Startup Week & task3 & 10 & Long-horizon latent-goal tracking under delayed effects. \\
    P2 & Incident Response Week & task6 & 7 & Shift adaptation under governance and recovery pressure. \\
    P3 & Quarterly Headcount Plan & task7 & 4 & Shorter planning-breadth task under reprioritization. \\
    \bottomrule
  \end{tabular}
  \caption{Paper-facing environments in the frozen benchmark. Internal ids are
  kept only for code and manifest provenance.}
  \label{tab:taskcatalog}
\end{table}

\section{Latent Objective and Public Schema}

The hidden objective is a structured latent state rather than a single class
label. It combines an archetype, a four-channel weight vector, a primary KPI,
and discrete operating attributes such as risk posture, planning horizon,
segment focus, and governance strictness. Table~\ref{tab:archetypes} gives the
base archetypes used in the frozen benchmark.

\begin{table}[!h]
  \centering
  \scriptsize
  \setlength{\tabcolsep}{5pt}
  \begin{tabular}{lp{0.24\linewidth}p{0.14\linewidth}p{0.36\linewidth}}
    \toprule
    Archetype & Channel weights & KPI & Operating attributes \\
    \midrule
    Growth & $(0.55, 0.15, 0.20, 0.10)$ & DAU & aggressive risk; quarterly planning; flexible governance; self-serve segment \\
    Retention & $(0.10, 0.60, 0.15, 0.15)$ & D30 retention & balanced risk; quarterly planning; strict governance; enterprise segment \\
    Revenue & $(0.10, 0.20, 0.55, 0.15)$ & MRR & balanced risk; quarterly planning; moderate governance; mid-market segment \\
    Efficiency & $(0.10, 0.15, 0.15, 0.60)$ & Ops margin & conservative risk; strategic planning; strict governance; enterprise segment \\
    \bottomrule
  \end{tabular}
  \caption{Base latent-goal archetypes. Channel order is growth, retention, revenue, efficiency.}
  \label{tab:archetypes}
\end{table}

In the frozen configuration, shifts occur with probability $0.6$; shift steps
are sampled from $\{3,4,5\}$; and the shift can be abrupt or drift-like. The
public schema exposes recurring customer accounts, stakeholders, internal
teams, market context, and governance constraints, while redacting the latent
channel decomposition and hidden utility weights.

\section{Formal Benchmark Definition}

\begin{table}[tbp]
  \centering
  \scriptsize
  \setlength{\tabcolsep}{5pt}
  \begin{tabular}{ll}
    \toprule
    Symbol & Meaning \\
    \midrule
    $x_t$ & hidden environment state at step $t$ \\
    $o_t$ & public observation exposed to the model \\
    $a_t$ & structured action emitted by the model \\
    $g_t$ & active hidden objective state \\
    $m_t$ & normalized four-channel public metric vector \\
    $w_t$ & active latent utility weight vector \\
    $U_t$ & latent utility under the active objective \\
    $q_t$ & normalized decision quality at step $t$ \\
    $\tilde{q}_t$ & stale-objective decision quality after a shift \\
    $A$ & adaptation score \\
    $C$ & trajectory coherence score \\
    $K$ & constraint-satisfaction score \\
    $\Delta_{\mathrm{shell}}$ & native-only minus overall score \\
    $\Delta_{\mathrm{shift}}$ & unshifted minus shifted score \\
    \bottomrule
  \end{tabular}
  \caption{Core notation used in the formal benchmark definition.}
  \label{tab:notation}
\end{table}

Let $x_t$ denote the hidden environment state at step $t$, let
$o_t = \psi(x_t)$ be the public observation exposed to the model, and let
$a_t \in \mathcal{A}$ be a structured action emitted through the benchmark's
action schema. An episode is a partially observed trajectory
\[
\tau = (o_0, a_0, o_1, a_1, \ldots, o_T).
\]
The hidden objective state is $g_t$, which determines which latent business
outcomes matter at time $t$. The agent never observes $g_t$ directly.

\paragraph{Latent utility.}
The environment maintains a normalized four-channel metric vector
$m_t \in [0,1]^4$ over growth, retention, revenue, and efficiency. The active
latent objective induces a weight vector $w_t \in \Delta^4$. The latent utility
used by the graders is
\[
U_t = \mathrm{clip}_{[0,1]} \left( \sum_{c \in \mathcal{C}} m_{t,c} \, w_{t,c} \right),
\qquad
\mathcal{C} = \{\text{growth}, \text{retention}, \text{revenue}, \text{efficiency}\}.
\]

When a hidden shift is present, the active objective changes from a base
objective $g^{(0)}$ to a shifted objective $g^{(1)}$ at step $t_\star$. The
implementation supports abrupt shifts and drift shifts. If the shift is
abrupt, $w_t = w^{(1)}$ for all $t \ge t_\star$. If the shift is a drift over
duration $d$, then the active weight vector is a normalized convex blend
\[
w_t = \mathrm{Normalize}\!\left((1-\rho_t)w^{(0)} + \rho_t w^{(1)}\right),
\qquad
\rho_t = \mathrm{clip}_{[0,1]}\!\left(\frac{t-t_\star+1}{d}\right).
\]
The shift is never explicitly announced to the model; it must be inferred from
changed consequences and public operating evidence.

For the three paper-facing environments, the overall score is a weighted combination of
final latent utility $U$, adaptation $A$, coherence $C$, and constraint
satisfaction $K$:
\[
S_{\mathrm{P1}} = 0.50\,U + 0.25\,A + 0.15\,C + 0.10\,K,
\]
\[
S_{\mathrm{P2}} = 0.4634\,U + 0.2927\,A + 0.1220\,C + 0.1220\,K,
\]
\[
S_{\mathrm{P3}} = 0.6364\,U + 0.2614\,A + 0.0568\,C + 0.0454\,K.
\]

\begin{table}[tbp]
  \centering
  \scriptsize
  \setlength{\tabcolsep}{5pt}
  \begin{tabular}{lrrrr}
    \toprule
    Task & Utility & Adaptation & Coherence & Constraints \\
    \midrule
    Startup Week & 0.5000 & 0.2500 & 0.1500 & 0.1000 \\
    Incident Response Week & 0.4634 & 0.2927 & 0.1220 & 0.1220 \\
    Quarterly Headcount Plan & 0.6364 & 0.2614 & 0.0568 & 0.0454 \\
    \bottomrule
  \end{tabular}
  \caption{Prespecified score coefficients for the paper-facing sequential
  tasks. These weights make adaptation and constraint behavior visible while
  preserving final latent utility as the largest single component.}
  \label{tab:scoreweights}
\end{table}

Let $q_t \in [0,1]$ be the normalized decision quality at step $t$, defined as
agent decision value divided by oracle decision value under the \emph{active}
objective. After a shift, the implementation also computes a stale-objective
quality $\tilde{q}_t$ whenever the base objective remains well-defined. The
first detected adaptation step for P1 and P2 is
\[
t_{\mathrm{det}} = \min \left\{ t \ge t_\star :
q_t \ge \max(0.72, \tilde{q}_t + 0.05) \right\},
\]
with the fallback rule $t_{\mathrm{det}} = T$ if no post-shift step crosses the
threshold. The remaining adaptation terms are
\[
R = \frac{1}{|D|} \sum_{t \in D} q_t,
\qquad
\mathrm{speed} = 1 - \frac{t_{\mathrm{det}} - t_\star}{\max(T - t_\star + 1, 1)},
\]
where $D = \{t_{\mathrm{det}}, \ldots, T\}$ if detection occurs and
$D=\varnothing$ otherwise. For P1 and P2, the adaptation score is
\[
A = 0.55 \cdot \mathrm{speed} + 0.45 \cdot R.
\]

P3 uses the same final weighted-score framework but a shorter-horizon
adaptation rule. Its detected adaptation step uses threshold
\[
q_t \ge \max(0.60, \tilde{q}_t + 0.02),
\]
with the same fallback of no post-detection recovery when the threshold is not
crossed. Let $R_7$ be mean post-detection decision quality, let
$\mathrm{speed}_7$ be the normalized detection-speed term over the post-shift
window, and let $C_{\mathrm{post}}$ be trajectory coherence over post-shift
actions. Then
\[
A_7 = (0.40 \cdot \mathrm{speed}_7 + 0.60 \cdot R_7) C_{\mathrm{post}}.
\]
When an episode has no hidden shift, the adaptation component is marked not
applicable and the weighted score is renormalized over the active components
only. When a shift exists but no post-shift decision row is available,
adaptation is scored as 0.0.

\begin{table}[tbp]
  \centering
  \scriptsize
  \setlength{\tabcolsep}{3pt}
  \begin{tabular}{p{0.11\linewidth}p{0.22\linewidth}p{0.19\linewidth}p{0.30\linewidth}}
    \toprule
    Task & Detection rule & Adaptation mix & Coherence and no-shift handling \\
    \midrule
    P1 Startup Week & $q_t \ge \max(0.72,\tilde{q}_t+0.05)$ & $0.55$ speed, $0.45$ recovery & Coherence ignores the true pivot; unshifted episodes omit adaptation and renormalize. \\
    P2 Incident Response Week & $q_t \ge \max(0.72,\tilde{q}_t+0.05)$ & $0.55$ speed, $0.45$ recovery & Coherence ignores the true pivot; unshifted episodes omit adaptation and renormalize. \\
    P3 Quarterly Headcount Plan & $q_t \ge \max(0.60,\tilde{q}_t+0.02)$ & $0.40$ speed, $0.60$ recovery & Adaptation is multiplied by post-shift coherence; unshifted episodes omit adaptation and renormalize. \\
    \bottomrule
  \end{tabular}
  \caption{Adaptation parameters for P1--P3. Shifted episodes without
  detection receive zero post-detection recovery; unshifted episodes omit
  adaptation and renormalize the score.}
  \label{tab:adaptparams}
\end{table}

Trajectory coherence is the mean within-regime cosine similarity between
successive action-strategy embeddings, ignoring the single pivot induced by a
true objective shift. To separate latent-policy failure from interface failure,
the paper also reports
\[
\Delta_{\mathrm{shell}} = S_{\mathrm{native}} - S_{\mathrm{overall}},
\qquad
\Delta_{\mathrm{shift}} = S_{\mathrm{unshifted}} - S_{\mathrm{shifted}},
\]
where $S_{\mathrm{native}}$ is the mean score over episodes that remain fully
native to the structured action interface and $S_{\mathrm{overall}}$ includes
fallback-affected episodes.

\section{Serving, Metadata, and Interface}

The strengthened paper protocol uses runtime version
\texttt{2026-04-07-shell-overhaul-v6} and freeze date \texttt{2026-04-18}. The
resolved manifest is the source of truth for the freeze: the analysis root
collects validated model runs, baseline runs, diagnostics, and targeted top-ups
under one dated paper bundle while preserving their original relative run
directories. The artifact bundle contains 10 locally served primary benchmark
configurations and one API-served external anchor across four served families:
\begin{itemize}[leftmargin=1.5em]
  \item Qwen3.5: 2B, 9B, 27B~\citep{qwen3};
  \item DeepSeek-R1: 1.5B, 8B, 14B~\citep{deepseekr1};
  \item Gemma4: 2B, 4B, 26B~\citep{gemma,gemma3};
  \item GPT-OSS: 20B locally, 120B as an external API anchor~\citep{gptoss}.
\end{itemize}

The strengthened bundle contains:
\begin{itemize}[leftmargin=1.5em]
  \item 630 primary model episode rows;
  \item 60 external-anchor episode rows;
  \item 360 baseline episode rows;
  \item 84 leakage-probe rows;
  \item generated tables, figures, case studies, and protocol descriptors tied
  to the frozen bundle.
\end{itemize}

The online environment also maintains a leakage-safe shaped reward used during
interaction, but not as the final benchmark metric.

The strengthened artifact bundle mixes two execution paths: local Ollama
serving for the 10 primary benchmark configurations and OpenRouter API serving
for GPT-OSS-120B. The provider difference is made explicit and GPT-OSS-120B is
reported only as an external anchor rather than as a load-bearing point in the
main scaling matrix.

\begin{table}[t]
  \centering
  \scriptsize
  \setlength{\tabcolsep}{4pt}
  \begin{tabular}{lllrrrrr}
    \toprule
    Model & Split & Provider & In tok (K) & Out tok (K) & Cost (\$) & Native & Empty fb \\
    \midrule
    DeepSeek 8B & core & Ollama local & 714.4 & 5.2 & 0.000 & 1.00 & 0.00 \\
    DeepSeek 8B & heldout & Ollama local & 717.8 & 5.4 & 0.000 & 1.00 & 0.00 \\
    DeepSeek 14B & core & Ollama local & 724.9 & 5.1 & 0.000 & 1.00 & 0.00 \\
    DeepSeek 14B & heldout & Ollama local & 706.2 & 5.3 & 0.000 & 1.00 & 0.00 \\
    Gemma 4B & core & Ollama local & 752.6 & 5.6 & 0.000 & 1.00 & 0.00 \\
    Gemma 4B & heldout & Ollama local & 747.2 & 5.6 & 0.000 & 1.00 & 0.00 \\
    Gemma 26B & core & Ollama local & 749.3 & 42.8 & 0.000 & 0.70 & 0.30 \\
    Gemma 26B & heldout & Ollama local & 748.2 & 38.8 & 0.000 & 0.73 & 0.27 \\
    GPT-OSS 20B & core & Ollama local & 812.5 & 16.8 & 0.000 & 1.00 & 0.00 \\
    GPT-OSS 20B & heldout & Ollama local & 784.2 & 8.6 & 0.000 & 1.00 & 0.00 \\
    Qwen 2B & core & Ollama local & 779.4 & 5.1 & 0.000 & 1.00 & 0.00 \\
    Qwen 2B & heldout & Ollama local & 785.1 & 5.0 & 0.000 & 1.00 & 0.00 \\
    Qwen 9B & core & Ollama local & 704.2 & 5.2 & 0.000 & 1.00 & 0.00 \\
    Qwen 9B & heldout & Ollama local & 694.2 & 5.8 & 0.000 & 1.00 & 0.00 \\
    Qwen 27B & core & Ollama local & 430.1 & 197.9 & 0.000 & 0.23 & 0.77 \\
    Qwen 27B & heldout & Ollama local & 430.1 & 196.9 & 0.000 & 0.47 & 0.53 \\
    \bottomrule
  \end{tabular}
  \caption{Provider/runtime summary for the frozen primary-task aggregate. This table is included because benchmark interpretation depends on serving path, token pressure, and interface reliability.}
  \label{tab:runtime}
\end{table}


The runtime table reports the visible serving path, token volume, cost, and
interface-reliability summaries for each primary model/split. Exact raw model
aliases, run roots, and role labels are preserved in
\texttt{metadata/paper\_protocol\_resolved.json} and
\texttt{metadata/run\_manifest.json}. The release also includes provider-path
and token-count records in \texttt{paper/release/COMPUTE.md}; the original
worker inventory is partially preserved, so the checklist marks full compute
inventory disclosure conservatively while keeping the runnable evaluation
artifact reproducible from the shipped manifest.

\section{Action Interface and Failure Taxonomy}

The action interface is intentionally narrow. On P1 and P2, a valid action
contains a task id plus a JSON-like payload with four core keys:
\texttt{chosen\_initiatives}, \texttt{messaging\_action},
\texttt{pricing\_change\_pct}, and \texttt{support\_policy}. Optional
rationale, memory, and belief notes may also be present. The allowed messaging
values are growth push, retention campaign, revenue upsell, cost comms, and
none. The allowed support policies are premium SLA, balanced triage,
automation first, and incident swarm. P3 uses the same action envelope but
replaces initiative/package selection with \texttt{budget\_allocations}.
Validation rejects irrelevant cross-task fields.

\begin{table}[tbp]
  \centering
  \scriptsize
  \setlength{\tabcolsep}{3pt}
  \begin{tabular}{lp{0.20\linewidth}p{0.16\linewidth}p{0.35\linewidth}}
    \toprule
    Signal & Operational meaning & Log field & Why it matters \\
    \midrule
    Native step & Action parses and satisfies the schema & strict action valid & Lets us attribute reward to the model's policy rather than fallback behavior. \\
    Empty fallback & Parser returns an empty/default action & empty fallback & Damaging in sequential tasks because the world advances while the model does little useful work. \\
    Cap hit & Response exceeds the output cap & response cap hit & Separates verbosity-related truncation from ordinary policy error. \\
    Parse repair & Action required bounded repair before execution & appendix diagnostic only & Tests whether failure is a lightweight parser issue or a deeper interface-control problem. \\
    Native episode rate & Fraction of episodes that remain native throughout & aggregated metric & Primary paper signal for shell reliability at the episode level. \\
    \bottomrule
  \end{tabular}
  \caption{Action-interface failure taxonomy used throughout the paper.}
  \label{tab:failtaxonomy}
\end{table}

\section{Supplementary Reference Tables}

The main model and baseline tables are reported in the main paper. This section
keeps only the external anchor and the diagnostics that are specific to the
supplementary analysis.

\begin{table}[t]
  \centering
  \scriptsize
  \setlength{\tabcolsep}{4pt}
  \begin{tabular}{lllrrrrr}
    \toprule
    Model & Split & Provider & Episodes & Overall & 95\% CI & Native rate & Empty fallback \\
    \midrule
    GPT-OSS 120B & core & OpenRouter API & 30 & 0.381 & [0.348, 0.415] & 0.60 & 0.40 \\
    GPT-OSS 120B & heldout & OpenRouter API & 30 & 0.365 & [0.339, 0.396] & 0.43 & 0.57 \\
    \bottomrule
  \end{tabular}
  \caption{External anchors excluded from the primary benchmark matrix because
  they were not run under the same local serving path. They are retained for
  transparency but are not used as load-bearing within-family scaling evidence.}
  \label{tab:externalanchors}
\end{table}


\begin{table}[t]
  \centering
  \small
  \begin{tabular}{p{0.24\linewidth}p{0.22\linewidth}cp{0.26\linewidth}}
    \toprule
    Check & Value & Passes & Threshold \\
    \midrule
    Core baseline ordering & 0.303 < 0.346 < 0.599 & Yes & random < visible heuristic < myopic oracle \\
    Core oracle gap & 0.296 & Yes & >= 0.20 \\
    Heldout baseline ordering & 0.310 < 0.371 < 0.598 & Yes & random < visible heuristic < myopic oracle \\
    Heldout oracle gap & 0.289 & Yes & >= 0.20 \\
    Silent-shift penalty & 0.172 & Yes & median unshifted-minus-shifted gap >= 0.10 \\
    \bottomrule
  \end{tabular}
  \caption{Selected benchmark health checks used in the paper.}
  \label{tab:health}
\end{table}


\begin{table}[t]
  \centering
  \scriptsize
  \setlength{\tabcolsep}{4pt}
  \begin{tabular}{llrrrrc}
    \toprule
    Variant & Split & Best primary & Best score & Oracle gap & Shift gap & Baseline order \\
    \midrule
    original & core & Qwen 27B & 0.443 & 0.156 & 0.154 & Yes \\
    original & heldout & DeepSeek 8B & 0.446 & 0.152 & 0.183 & Yes \\
    utility-only & core & Qwen 27B & 0.392 & 0.012 & 0.072 & Yes \\
    utility-only & heldout & GPT-OSS 20B & 0.393 & 0.016 & 0.083 & Yes \\
    balanced & core & DeepSeek 14B & 0.494 & 0.087 & 0.185 & Yes \\
    balanced & heldout & DeepSeek 8B & 0.499 & 0.081 & 0.209 & Yes \\
    adaptation-heavy & core & Qwen 27B & 0.463 & 0.089 & 0.214 & Yes \\
    adaptation-heavy & heldout & DeepSeek 8B & 0.462 & 0.088 & 0.249 & Yes \\
    adaptation-light & core & DeepSeek 1.5B & 0.519 & 0.057 & 0.109 & Yes \\
    adaptation-light & heldout & Qwen 2B & 0.523 & 0.056 & 0.118 & Yes \\
    no-coherence & core & Qwen 27B & 0.430 & 0.061 & 0.169 & No \\
    no-coherence & heldout & DeepSeek 8B & 0.429 & 0.060 & 0.207 & Yes \\
    \bottomrule
  \end{tabular}
  \caption{Offline score-weight sensitivity from frozen final-episode grader
  outputs. ``Shift gap'' is median unshifted-minus-shifted gap across primary
  models. The utility-only variant is included to show that the benchmark's
  main signal depends on the intended adaptation construct, not only final
  utility.}
  \label{tab:scoresensitivity}
\end{table}


The sensitivity table is best read as a construct check. Baseline ordering is
stable for the original, utility-only, balanced, adaptation-heavy, and
adaptation-light variants on both splits. The no-coherence variant is included
as a useful boundary case: it removes one piece of the sequential-behavior
signal and breaks core baseline ordering. This supports keeping coherence as a
small but nonzero component of the prespecified score. The utility-only rows
also clarify why LatentGoalOps should not be interpreted as a final-state
utility benchmark alone: once adaptation is removed, the myopic-oracle gap
shrinks sharply even though the shifted-episode penalty remains visible.

\begin{table}[t]
  \centering
  \scriptsize
  \setlength{\tabcolsep}{4pt}
  \begin{tabular}{llrrrrrr}
    \toprule
    Task & Split & Full & Numeric & Sender & Order & Max restricted & Goal tokens \\
    \midrule
    Startup Week & core & 0.20 & 0.05 & 0.25 & 0.25 & 0.30 & 0.00 \\
    Startup Week & heldout & 0.20 & 0.15 & 0.20 & 0.25 & 0.25 & 0.00 \\
    Incident Response Week & core & 0.20 & 0.10 & 0.30 & 0.25 & 0.30 & 0.00 \\
    Incident Response Week & heldout & 0.15 & 0.20 & 0.30 & 0.25 & 0.30 & 0.00 \\
    Quarterly Headcount Plan & core & 0.25 & 0.25 & 0.25 & 0.25 & 0.35 & 0.00 \\
    Quarterly Headcount Plan & heldout & 0.20 & 0.15 & 0.25 & 0.25 & 0.30 & 0.00 \\
    \bottomrule
  \end{tabular}
  \caption{Compact restricted-view leakage summary for paper-facing tasks.
  Chance is 0.25 for the four-way archetype probe. ``Max restricted'' is the
  maximum over non-full views in the frozen audit and is interpreted as a
  screening diagnostic rather than a proof of no shortcut.}
  \label{tab:leakagecompact}
\end{table}


\begin{table}[t]
  \centering
  \scriptsize
  \setlength{\tabcolsep}{4pt}
  \begin{tabular}{lrrrrrrrr}
    \toprule
    \multirow{2}{*}{Model} & \multicolumn{2}{c}{Core P1} & \multicolumn{2}{c}{Core P2} & \multicolumn{2}{c}{Heldout P1} & \multicolumn{2}{c}{Heldout P2} \\
    & Adapt. & Zero rate & Adapt. & Zero rate & Adapt. & Zero rate & Adapt. & Zero rate \\
    \midrule
    Qwen 2B & 0.000 & 1.00 & 0.000 & 1.00 & 0.000 & 1.00 & 0.000 & 1.00 \\
    Qwen 9B & 0.000 & 1.00 & 0.181 & 0.78 & 0.198 & 0.67 & 0.000 & 1.00 \\
    Qwen 27B & 0.117 & 0.82 & 0.060 & 0.91 & 0.000 & 1.00 & 0.000 & 1.00 \\
    DeepSeek 1.5B & 0.000 & 1.00 & 0.000 & 1.00 & 0.000 & 1.00 & 0.000 & 1.00 \\
    DeepSeek 8B & 0.101 & 0.89 & 0.190 & 0.78 & 0.180 & 0.78 & 0.185 & 0.78 \\
    DeepSeek 14B & 0.090 & 0.89 & 0.195 & 0.78 & 0.195 & 0.67 & 0.000 & 1.00 \\
    Gemma 2B & 0.000 & 1.00 & 0.000 & 1.00 & 0.000 & 1.00 & 0.000 & 1.00 \\
    Gemma 4B & 0.132 & 0.78 & 0.181 & 0.78 & 0.237 & 0.67 & 0.000 & 1.00 \\
    Gemma 26B & 0.161 & 0.78 & 0.298 & 0.61 & 0.000 & 1.00 & 0.033 & 0.94 \\
    GPT-OSS 20B & 0.098 & 0.78 & 0.277 & 0.67 & 0.057 & 0.89 & 0.000 & 1.00 \\
    \bottomrule
  \end{tabular}
  \caption{Supplementary adaptation audit. ``Adapt.'' is mean adaptation score
  on shifted episodes; ``Zero rate'' is the fraction of shifted episodes with
  zero adaptation score.}
  \label{tab:adaptationaudit}
\end{table}


\begin{table}[t]
  \centering
  \scriptsize
  \setlength{\tabcolsep}{4pt}
  \begin{tabular}{llrrrrr}
    \toprule
    Model & Split & Frozen overall & Shell overall & Shell assisted & Native rate & Repair rate \\
    \midrule
    Gemma 26B & core & 0.395 & 0.285 & 0.278 & 0.00 & 0.10 \\
    Gemma 26B & heldout & 0.361 & 0.287 & 0.353 & 0.00 & 0.10 \\
    Qwen 27B & core & 0.361 & 0.333 & 0.340 & 0.05 & 0.40 \\
    Qwen 27B & heldout & 0.385 & 0.345 & 0.375 & 0.25 & 0.20 \\
    \bottomrule
  \end{tabular}
  \caption{Appendix shell-controlled diagnostic with parse repair enabled and heuristic rescue disabled. ``Shell assisted'' is the assisted aggregate reported by the diagnostic runner.}
  \label{tab:shelldiagnostic}
\end{table}


The shell-controlled diagnostic is an appendix mechanism check focused on the
two week-long anchor tasks, Startup Week and Incident Response Week. It is not
used as a replacement for the headline three-task aggregate; instead, it tests
whether parse repair alone explains the largest interface-mediated drops.

\begin{table}[t]
  \centering
  \tiny
  \setlength{\tabcolsep}{3pt}
  \begin{tabular}{llp{0.23\linewidth}rrrp{0.18\linewidth}}
    \toprule
    Family & Split & Transition & $\Delta$ overall & $\Delta$ native-only & $\Delta$ native rate & Diagnosis \\
    \midrule
    DeepSeek-R1 & core & DeepSeek 1.5B $\rightarrow$ DeepSeek 8B & -0.026 & -0.026 & 0.00 & native policy regression \\
    DeepSeek-R1 & heldout & DeepSeek 1.5B $\rightarrow$ DeepSeek 8B & 0.016 & 0.016 & 0.00 & gain \\
    DeepSeek-R1 & core & DeepSeek 8B $\rightarrow$ DeepSeek 14B & 0.035 & 0.035 & 0.00 & gain \\
    DeepSeek-R1 & heldout & DeepSeek 8B $\rightarrow$ DeepSeek 14B & -0.017 & -0.017 & 0.00 & native policy regression \\
    Gemma4 & core & Gemma 2B $\rightarrow$ Gemma 4B & 0.016 & 0.016 & 0.00 & gain \\
    Gemma4 & heldout & Gemma 2B $\rightarrow$ Gemma 4B & 0.013 & 0.013 & 0.00 & gain \\
    Gemma4 & core & Gemma 4B $\rightarrow$ Gemma 26B & -0.025 & 0.013 & -0.30 & interface regression \\
    Gemma4 & heldout & Gemma 4B $\rightarrow$ Gemma 26B & -0.067 & -0.044 & -0.27 & interface regression \\
    Qwen3.5 & core & Qwen 2B $\rightarrow$ Qwen 9B & -0.008 & -0.008 & 0.00 & flat \\
    Qwen3.5 & heldout & Qwen 2B $\rightarrow$ Qwen 9B & -0.013 & -0.013 & 0.00 & native policy regression \\
    Qwen3.5 & core & Qwen 9B $\rightarrow$ Qwen 27B & 0.021 & 0.022 & -0.03 & gain \\
    Qwen3.5 & heldout & Qwen 9B $\rightarrow$ Qwen 27B & -0.037 & 0.035 & -0.53 & interface regression \\
    \bottomrule
  \end{tabular}
  \caption{Scaling-transition summary. Native-only deltas help separate clean policy regressions from interface-mediated regressions.}
  \label{tab:scalingtransitions}
\end{table}


Full per-environment and subscore tables are shipped as machine-readable CSV
files and LaTeX fragments under \texttt{outputs/paper-freeze-2026-04-18/tables}
and \texttt{paper/latex/generated}. They are intentionally not expanded here so
the supplement remains readable and free of float-only table pages.

\section{Supplementary Figures}

\begin{figure}[tbp]
  \centering
  \includegraphics[width=0.92\linewidth]{\figdir/overall_vs_native.pdf}
\vspace{0.3em}
  \includegraphics[width=0.92\linewidth]{\figdir/action_reliability.pdf}
  \caption{Primary interface diagnostics. Top: observed score versus
  native-only score. Bottom: native-action, empty-fallback, and cap-hit rates
  on the frozen primary-task aggregate.}
  \label{fig:supp:native-reliability}
\end{figure}

\begin{figure}[tbp]
  \centering
  \includegraphics[width=0.92\linewidth]{\figdir/shifted_vs_unshifted.pdf}
  \caption{Shifted episodes remain materially harder than no-shift episodes on
  the frozen primary-task aggregate. Error bars denote 95\% bootstrap
  confidence intervals for the score aggregates.}
  \label{fig:supp:shifted}
\end{figure}

\begin{figure}[tbp]
  \centering
  \includegraphics[width=0.92\linewidth]{\figdir/task_heatmap.pdf}
\vspace{0.3em}
  \includegraphics[width=0.92\linewidth]{\figdir/subscore_heatmap.pdf}
  \caption{Primary-task heatmaps. Top: task-level performance on the frozen
  main-task core. Bottom: subscore decomposition across utility, adaptation,
  coherence, and constraints.}
  \label{fig:supp:heatmaps}
\end{figure}

\begin{figure}[tbp]
  \centering
  \includegraphics[width=0.92\linewidth]{\figdir/shell_controlled_diagnostic.pdf}
\vspace{0.3em}
  \includegraphics[width=0.92\linewidth]{\figdir/token_pressure_vs_native_rate.pdf}
  \caption{Interface-mechanism diagnostics. Top: shell-controlled diagnostic;
  parse repair produces only modest improvement for Qwen 27B and does not
  rescue Gemma 26B. Bottom: token pressure versus native-action reliability;
  the most shell-fragile large-model runs are also the most verbose.}
  \label{fig:supp:shell-token}
\end{figure}

\begin{figure}[tbp]
  \centering
  \includegraphics[width=0.92\linewidth]{\figdir/adaptation_audit.pdf}
\vspace{0.3em}
  \includegraphics[width=0.92\linewidth]{\figdir/leakage_audit.pdf}
  \caption{Adaptation and leakage diagnostics. Top: adaptation audit showing
  mean adaptation score and zero-adaptation episode rate by task and split.
  Bottom: leakage audit across restricted observation views on the primary
  aggregate tasks. Chance is 0.25 for the four-way archetype probe, and the
  explicit-goal-token rate is 0.0 throughout the frozen audit.}
  \label{fig:supp:adapt-leakage}
\end{figure}

\begin{figure}[tbp]
  \centering
  \includegraphics[width=0.92\linewidth]{\figdir/heldout_gap.pdf}
\vspace{0.3em}
  \includegraphics[width=0.92\linewidth]{\figdir/reliability_ci.pdf}
  \caption{Generalization and repeatability diagnostics. Top: heldout-minus-core
  gap on the primary aggregate tasks. Bottom: repeat-based reliability
  estimates where repeat coverage exists; GPT-OSS 120B remains a low-coverage
  appendix anchor rather than a load-bearing comparison.}
  \label{fig:supp:transfer-reliability}
\end{figure}

\subsection*{Leakage Audit Details}

The leakage figure condenses the frozen leakage table
\texttt{tables/leakage\_results.csv} inside the paper-freeze artifact bundle.
On the primary tasks, the full-observation centroid probe is 0.20 on core
Startup Week, 0.20 on core Incident Response Week, 0.25 on core Quarterly
Headcount Plan, and 0.15--0.20 on the corresponding heldout slices.
Restricted views such as \texttt{numeric\_only},
\texttt{sender\_only}, and \texttt{stakeholder\_order\_only} remain near the
four-way chance regime, while \texttt{explicit\_goal\_token\_rate} is 0.0 in
every frozen probe output.

\section{Case Studies and Provenance}

The frozen bundle includes four qualitative trajectory analyses:
\begin{itemize}[leftmargin=1.5em]
  \item \texttt{deepseek14\_success.md}
  \item \texttt{qwen27\_interface\_failure.md}
  \item \texttt{gemma26\_interface\_failure.md}
  \item \texttt{qwen9\_native\_policy\_regression.md}
\end{itemize}

\begin{table}[H]
  \centering
  \scriptsize
  \setlength{\tabcolsep}{5pt}
  \begin{tabular}{lp{0.13\linewidth}p{0.18\linewidth}cp{0.34\linewidth}}
    \toprule
    Model & Split & Task & Score & Dominant diagnosis \\
    \midrule
    DeepSeek 14B & core & Incident Response Week & 0.6353 & Fully native, coherent incident-focused policy that keeps support and messaging aligned while delayed effects land. \\
    Qwen 27B & core & Incident Response Week & 0.2564 & Starts with a plausible action, then falls into repeated empty fallbacks after the hidden shift. \\
    Gemma 26B & heldout & Startup Week & 0.2561 & Mixed shell/policy failure with repeated empty steps, cap hits, and almost no recovery traction. \\
    Qwen 9B & core & Startup Week & 0.3070 & Fully native trajectory that stays schema-correct but keeps reverting to stale startup moves after the hidden shift. \\
    \bottomrule
  \end{tabular}
  \caption{Curated case-study summary used in the supplement discussion.}
  \label{tab:casestudies}
\end{table}

Table~\ref{tab:trajectory-excerpts} adds transcript-style trajectory snippets.
Unlike a dialogue benchmark, LatentGoalOps exposes structured operating state
and structured actions, so each row pairs a compact state snapshot with a
paraphrase of the executed action rather than dumping raw JSON payloads.

Across the four frozen case studies, similar scalar scores can arise from
meaningfully different mechanisms. DeepSeek 14B provides the clean positive
control: the action shell remains native and the policy stays incident-focused
while delayed effects resolve. Qwen 27B shows a different pattern, where one
mid-episode shell break is followed by a policy over-correction into aggressive
startup moves. Qwen 9B is useful because it removes shell collapse from the
story almost entirely: the trajectory stays schema-correct but keeps reverting
to stale growth behavior after the hidden shift. Gemma 26B combines both
failure classes, with repeated empty actions and almost no policy recovery.
That qualitative separation is why the paper reports strict score, native-only
score, and native-episode rate together rather than collapsing everything into
a single scalar.

\begin{table*}[tbp]
  \centering
  \footnotesize
  \setlength{\tabcolsep}{4pt}
  \begin{tabular}{p{0.08\textwidth}p{0.15\textwidth}p{0.26\textwidth}p{0.38\textwidth}}
    \toprule
    Step/day & Observed state & Executed action (paraphrased) & What the snippet shows \\
    \midrule
    \multicolumn{4}{l}{\textbf{DeepSeek 14B, core Incident Response Week, score 0.6353, fully native}} \\
    Step 0 / Day 2 & 3 pending effects; early incident week state still unstable & Cut usage pricing 10\%, keep retention messaging, keep incident-swarm support, push usage-pricing initiative & Starts with a service-preserving move even before delayed effects land; no shell instability and no thrashing across goals. \\
    Step 2 / Day 4 & 4 pending effects; incident backlog still active & Optimize infra while keeping retention messaging and incident-swarm support & Mid-episode policy stays incident-focused instead of pivoting to generic growth or cost moves. \\
    Step 5 / Day 7 & 2 pending effects; delayed effects still resolving & Refactor infra, keep retention messaging, keep incident-swarm support, cut usage pricing 10\% & Strong positive reward arrives without the model abandoning its earlier operational stance. \\
    Step 6 / Day 8 & 0 pending effects; week closes cleanly & Refactor incident tooling, keep retention messaging, keep incident-swarm support & Coherent finish: action shell stays native, and policy remains aligned through episode end. \\
    \midrule
    \multicolumn{4}{l}{\textbf{Qwen 27B, core Startup Week, score 0.3685, shell collapse after hidden shift}} \\
    Step 0 / Day 2 & 5 pending effects; early startup state & Optimize infra, refactor incident tooling, automate support triage; cost-oriented messaging; automation-first support & Plausible pre-shift startup move; qualitative issue is not immediate nonsense but what happens after state changes. \\
    Step 3 / Day 5 & 0 pending effects; post-shift transition point & Empty fallback action \{\} after parse failure & Critical shell break: world advances, but model contributes nothing useful at exactly the point where adaptation pressure matters most. \\
    Step 6 / Day 8 & 5 pending effects; recovery attempt under new state & Ship usage pricing, switch to revenue upsell, move to premium-SLA support & Over-corrects into aggressive revenue posture after the shell break instead of recovering a stable startup policy. \\
    Steps 7--9 / Days 9--11 & Pending effects shrink, but score remains negative & Repeats referral-loop plus growth-push moves despite negative rewards & Failure mode is not only one malformed step; policy stays stale and misaligned even after shell recovers. \\
    \bottomrule
  \end{tabular}
  \caption{Transcript-style trajectory excerpts paraphrased from the frozen
  case-study markdowns. These are qualitative observation-action snapshots,
  not chain-of-thought traces. They make two distinct failure modes concrete:
  coherent native adaptation versus interface collapse followed by policy
  over-correction.}
  \label{tab:trajectory-excerpts}
\end{table*}

The strengthened bundle is organized to support direct regeneration of the
paper assets:

\begin{verbatim}
outputs/paper-freeze-2026-04-18/
  tables/
  figures/
  case_studies/
  metadata/
\end{verbatim}

The \texttt{tables/} and \texttt{figures/} directories contain the analysis
outputs consumed by the manuscript. The \texttt{case\_studies/} directory
stores curated trajectory narratives, and \texttt{metadata/} contains the
frozen manifest and protocol descriptors. The exact upstream raw run
directories needed to rebuild those assets are enumerated in
\texttt{metadata/run\_manifest.json} and are shipped in the anonymous reviewer
archive under their original relative paths rather than duplicated inside the
\texttt{2026-04-18} analysis root itself.

\section{Reproducibility}

The paper assets are generated with:

\begin{verbatim}
uv run paper-assets --config configs/paper_protocol_strengthened.yaml
\end{verbatim}

This command rebuilds the paper tables, figures, and case studies from the
strengthened artifact bundle. The manuscript itself can be compiled with the
bundled \texttt{tectonic} binary included in the repository tools directory.
The submission's anonymous review archive contains the code, the
\texttt{outputs/paper-freeze-2026-04-18} analysis root, the exact raw run
directories referenced by the manifest, and the release documents
\texttt{paper/release/REVIEW\_BUNDLE.md}, \texttt{paper/release/COMPUTE.md},
and \texttt{paper/release/THIRD\_PARTY\_ASSETS.md}. Those materials make the
benchmark runnable for reviewers while keeping the main paper focused on the
benchmark claims rather than on packaging detail.

\section{Threats to Validity}

\subsection{Provider heterogeneity}

Most frozen runs are local Ollama runs, while GPT-OSS-120B is served through
OpenRouter. We therefore treat GPT-OSS-120B as an external anchor and not as a
clean within-family local scaling point.

\subsection{Interface confounding}

Some large-model drops are clearly interface-mediated. Rather than hiding that
fact, the paper makes it part of the result and reports strict overall score,
native-only score, fallback rates, and the shell-controlled diagnostic.

\subsection{Adaptation metric sharpness}

The adaptation score for P1 and P2 uses a thresholded behavioral detection
rule. That choice is operationally meaningful, but it also makes some slices of
the benchmark harsher than a fully smooth scalar metric would. The supplement
therefore includes both the formal definition and a dedicated adaptation audit.

\subsection{Synthetic world construction}

LatentGoalOps is synthetic by design. The claim is not that it replicates every
real operating environment; it is that it isolates a capability that existing
benchmarks often blur: hidden-objective tracking with delayed effects and
silent non-stationarity. The construct-validity argument therefore relies on
heldout splits, leakage audits, task design transparency, and baseline
separation rather than on realism language alone.

\fi

\ifincludechecklist
  \clearpage
  \section*{NeurIPS Paper Checklist}

\begin{enumerate}

\item {\bf Claims}
    \item[] Question: Do the main claims made in the abstract and introduction accurately reflect the paper's contributions and scope?
    \item[] Answer: \answerYes{}
    \item[] Justification: The abstract and introduction state the benchmark contribution, the latent-adaptation scope, the heterogeneous scaling pattern, and the interface-reliability finding in a way that matches the reported results and stated caveats; see Abstract, Section 1, and Section 5.

\item {\bf Limitations}
    \item[] Question: Does the paper discuss the limitations of the work performed by the authors?
    \item[] Answer: \answerYes{}
    \item[] Justification: The main paper contains a dedicated Limitations section covering provider heterogeneity, interface confounding, adaptation-metric sharpness, and synthetic-world scope; see Section 6.

\item {\bf Theory assumptions and proofs}
    \item[] Question: For each theoretical result, does the paper provide the full set of assumptions and a complete (and correct) proof?
    \item[] Answer: \answerNA{}
    \item[] Justification: The paper is an empirical benchmark paper and does not present formal theorems or proofs. The equations define benchmark mechanics and grading, not new theoretical results.

\item {\bf Experimental result reproducibility}
    \item[] Question: Does the paper fully disclose all the information needed to reproduce the main experimental results of the paper to the extent that it affects the main claims and/or conclusions of the paper (regardless of whether the code and data are provided or not)?
    \item[] Answer: \answerYes{}
    \item[] Justification: The manuscript and release bundle specify the frozen protocol, task scope, seed counts, split structure, model scope, regeneration commands, resolved manifest, and artifact layout; exploratory hosted API probes are labeled separately and their raw logs are included under \texttt{outputs/api-probe-2026-05-07}. See Sections 2--4 and the release files \texttt{paper/release/README.md} and \texttt{paper/release/REPRODUCTION.md}.

\item {\bf Open access to data and code}
    \item[] Question: Does the paper provide open access to the data and code, with sufficient instructions to faithfully reproduce the main experimental results, as described in supplemental material?
    \item[] Answer: \answerYes{}
    \item[] Justification: The submission now includes an anonymous reviewer archive that packages the runnable benchmark code, the frozen \texttt{2026-04-18} analysis root, the exact raw run directories referenced by \texttt{metadata/run\_manifest.json}, the exploratory hosted API probe logs, and reproduction instructions; see Section 4 and the release files \texttt{paper/release/README.md}, \texttt{paper/release/REPRODUCTION.md}, and \texttt{paper/release/REVIEW\_BUNDLE.md}.

\item {\bf Experimental setting/details}
    \item[] Question: Does the paper specify all the training and test details (e.g., data splits, hyperparameters, how they were chosen, type of optimizer) necessary to understand the results?
    \item[] Answer: \answerYes{}
    \item[] Justification: The paper and release package specify the benchmark tasks, hidden-shift configuration, split structure, seed counts, scoring definitions, serving path, and model roles needed to interpret the evaluation; hosted API probes additionally disclose provider route, model identifiers, split coverage, and their exploratory status. See Sections 2--4 and the release protocol files.

\item {\bf Experiment statistical significance}
    \item[] Question: Does the paper report error bars suitably and correctly defined or other appropriate information about the statistical significance of the experiments?
    \item[] Answer: \answerYes{}
    \item[] Justification: The main result and baseline tables report 95\% bootstrap confidence intervals on the headline aggregates, and the text explains the role of native-only estimates and effective sample size; see Section 4 and the generated result tables.

\item {\bf Experiments compute resources}
    \item[] Question: For each experiment, does the paper provide sufficient information on the computer resources (type of compute workers, memory, time of execution) needed to reproduce the experiments?
    \item[] Answer: \answerNo{}
    \item[] Justification: The release package reports provider path, token counts, preserved wall-clock timing, and a validation-host note in \texttt{paper/release/COMPUTE.md}, but the exact original worker inventory for the full frozen runs was not preserved in the manifest. We therefore keep the conservative No answer rather than over-claim hardware-normalized reproducibility.

\item {\bf Code of ethics}
    \item[] Question: Does the research conducted in the paper conform, in every respect, with the NeurIPS Code of Ethics \url{https://neurips.cc/public/EthicsGuidelines}?
    \item[] Answer: \answerYes{}
    \item[] Justification: To the best of our knowledge, yes. The benchmark is synthetic, does not involve human subjects, and the paper and datasheet explicitly restrict inappropriate deployment-style uses; see Section 6 and \texttt{paper/release/DATASHEET.md}.

\item {\bf Broader impacts}
    \item[] Question: Does the paper discuss both potential positive societal impacts and negative societal impacts of the work performed?
    \item[] Answer: \answerYes{}
    \item[] Justification: The release datasheet and the paper's limitations discussion address appropriate uses, inappropriate uses, and the risk of over-interpreting a synthetic benchmark as a real deployment simulator; see Section 6 and \texttt{paper/release/DATASHEET.md}.

\item {\bf Safeguards}
    \item[] Question: Does the paper describe safeguards that have been put in place for responsible release of data or models that have a high risk for misuse (e.g., pre-trained language models, image generators, or scraped datasets)?
    \item[] Answer: \answerNA{}
    \item[] Justification: The paper releases a synthetic benchmark and frozen evaluation artifacts, not a high-risk generative model or scraped dataset. The benchmark also includes intended-use and limitation guidance in the release datasheet.

\item {\bf Licenses for existing assets}
    \item[] Question: Are the creators or original owners of assets (e.g., code, data, models), used in the paper, properly credited and are the license and terms of use explicitly mentioned and properly respected?
    \item[] Answer: \answerYes{}
    \item[] Justification: The release package now includes a consolidated asset-and-terms inventory for the evaluated external model families and serving dependencies, with official source links and notes on local alias handling; see \texttt{paper/release/THIRD\_PARTY\_ASSETS.md}.

\item {\bf New assets}
    \item[] Question: Are new assets introduced in the paper well documented and is the documentation provided alongside the assets?
    \item[] Answer: \answerYes{}
    \item[] Justification: The benchmark release includes documented code, a frozen artifact bundle, reproduction instructions, a datasheet, and Croissant metadata; see Section 4 and the files under \texttt{paper/release/}.

\item {\bf Crowdsourcing and research with human subjects}
    \item[] Question: For crowdsourcing experiments and research with human subjects, does the paper include the full text of instructions given to participants and screenshots, if applicable, as well as details about compensation (if any)?
    \item[] Answer: \answerNA{}
    \item[] Justification: The paper does not involve crowdsourcing or human-subject experiments.

\item {\bf Institutional review board (IRB) approvals or equivalent for research with human subjects}
    \item[] Question: Does the paper describe potential risks incurred by study participants, whether such risks were disclosed to the subjects, and whether Institutional Review Board (IRB) approvals (or an equivalent approval/review based on the requirements of your country or institution) were obtained?
    \item[] Answer: \answerNA{}
    \item[] Justification: The paper does not involve crowdsourcing or human-subject experiments, so IRB approval is not applicable.

\item {\bf Declaration of LLM usage}
    \item[] Question: Does the paper describe the usage of LLMs if it is an important, original, or non-standard component of the core methods in this research? Note that if the LLM is used only for writing, editing, or formatting purposes and does \emph{not} impact the core methodology, scientific rigor, or originality of the research, declaration is not required.
    \item[] Answer: \answerYes{}
    \item[] Justification: LLMs are the benchmarked systems in the experimental evaluation, including locally served models and hosted OpenRouter probes. The benchmark construction, deterministic scoring, and analysis pipeline do not rely on an auxiliary LLM-based labeling or generation procedure.

\end{enumerate}

\fi

\end{document}
